\subsection{Complexity ledger and the next certification boundary}\label{sec:mv:cost}

\begin{definition}[Compiled proof counts]\label{def:mv:proof-count}
Let $B=sn$ be the number of registered
binary commitments for $s$ field
elements and let $Q\le nE$ count the
nonempty affine-parity rows for
$E$ field-affine equations.
Each bit commitment has one bitness
OR statement and each nonempty row
has one aggregate OR statement.
Write $J:=B+Q$ for this number
of prescribed OR statements.
Each statement is replaced by
$\kappa$ sequential mini-proofs.
\end{definition}

\begin{lemma}[Exact mini-proof count]\label{lem:mv:proof-count}
The modified compiler performs
\[
 \boxed{M=\kappa J=\kappa(B+Q)}
\]
one-bit OR mini-proofs, provided
none is skipped by early rejection.
Its verifier sends $M$ challenge
bits, one after each OR first
message, in addition to the
outer field challenges.
\end{lemma}
\begin{proof}
There are $B$ bitness statements
and $Q$ aggregate statements
by Definition~\ref{def:mv:proof-count}.
Each runs exactly $\kappa$ copies
of Definition~\ref{def:mv:mini}
under a non-rejecting execution.
Thus their sum is
$\kappa(B+Q)$.
Each mini-proof receives exactly
one bit challenge after sending
its first message.
\end{proof}

\begin{lemma}[Expected and bounded simulator costs]\label{lem:mv:simulation-cost}
For the $M$ mini-proofs of
Lemma~\ref{lem:mv:proof-count},
the exact malicious-verifier
simulator uses at most $2M$
expected mini-proof trials.
A truncated simulator with
\[
 R=\lambda+\lceil\log_2 M\rceil
\]
uses at most $MR$ trials and
has statistical simulation
error at most $2^{-\lambda}$.
\end{lemma}
\begin{proof}
Apply
Lemma~\ref{lem:mv:repeat-cost}
to all $M$ possible copies
to obtain the expectation.
Apply
Theorem~\ref{thm:mv:truncated}
with the displayed $R$
for the strict upper bound
and its error estimate.
\end{proof}

\begin{proposition}[A security-efficiency tradeoff]\label{prop:mv:tradeoff}
Amplifying the inner challenge
with $\kappa$ \emph{sequential}
one-bit repetitions reduces the
baseline challenge-guessing
knowledge error to $2^{-\kappa}$
while increasing raw proof
communication by a factor of
$\kappa$. The resulting
uncompressed affine compiler
remains nonsuccinct.
\end{proposition}
\begin{proof}
A deterministic prover with no
two accepting challenge strings
passes at most one of the
$2^\kappa$ challenge sequences,
by Lemma~\ref{lem:mv:single-leaf}.
Thus its baseline acceptance is
at most $2^{-\kappa}$.
Lemma~\ref{lem:mv:proof-count}
shows that each prescribed OR
statement is transmitted
$\kappa$ times, using fresh
first messages and responses.
The size per copy remains
proportional to its branch set
by Lemma~\ref{lem:realize:or-communication},
so the raw communication is
multiplied by $\kappa$.
The direct compiler was already
nonsuccinct by
Proposition~\ref{prop:realize:not-succinct};
sequential repetition does not
add a compression step.
\end{proof}

\begin{definition}[Lean certification graph]\label{def:mv:lean-graph}
The intended dependency order
of Chapter~11 is:
\begin{enumerate}[label=\textup{(\roman*)}]
\item fixed-bit OR uniform
hyperplane and conditional
transcript distribution;
\item balanced malicious-verifier
simulator and abort-preserving
selection lemma;
\item sequential transcript
composition and bounded retries;
\item challenge-tree forking,
same-first-message extraction
and the quantitative $\varepsilon$
inequality;
\item stand-alone temporal
extraction and false-parity
discrete-log reduction;
\item a separately constructed
global temporal extractor for
all dependent commitments;
\item adaptive KBFold oracle
opening simulation and
post-quantum soundness.
\end{enumerate}
Items (vi)--(vii) are
\emph{unproved obligations},
not Lean declarations claimed
to exist.
\end{definition}

\begin{remark}[Historical placement]
The one-bit challenge rewinding
strategy is classical
zero-knowledge methodology
\cite{Goldreich2001Foundations};
our concrete OR transcript
uses the established Schnorr and
Cramer--Damg{\aa}rd--Schoenmakers
constructions
\cite{Schnorr1991,CramerDamgardSchoenmakers1994}.
The mathematical work here is to
make the challenge range, abort
conditioning, sequential
composition, exact forking
probability and extraction
timing explicit for the
kernel-coordinate affine
Ring Switching interface.
No claim that one-bit rewinding
itself is new is made.
\end{remark}

\begin{remark}[Next substantive mathematical problem]
To close the Chapter~10 early
extraction obligation, the
multistage extractor must produce
\emph{one} fixed hidden field
tuple before each outer challenge,
with the same accepted public
history and independent outer
challenge law. Only after that
can the resulting malicious-verifier
ZK affine proof be safely
composed with KBFold's original
authenticated folding-oracle
transcript, and its odd-
characteristic proximity theorem.
\end{remark}
